\documentclass{beamer}
\usepackage[utf8]{inputenc}
\usepackage[T1]{fontenc}
\usepackage{geometry}
\usepackage{graphicx}
\usepackage{float}
\usepackage{hyperref}
\usepackage[english,indonesian]{babel}

\title{Optimasi Pemrosesan Template LaTeX Menggunakan Pendekatan Berbasis AI}
\author{Academic Support Agent}
\date{\today}

\begin{document}
\maketitle

\begin{abstract}

\end{abstract}

\chapter{PENDAHULUAN}


\chapter{TINJAUAN PUSTAKA}


% Original content (commented out):
% \documentclass[10pt,aspectratio=1610,compress,dvipsnames]{beamer}

\usetheme[
%%% option passed to the outer theme
%    progressstyle=fixedCircCnt,   % fixedCircCnt, movingCircCnt (moving is deault)
  ]{Feather}
  
% If you want to change the colors of the various elements in the theme, edit and uncomment the following lines

% Change the bar colors:
%\setbeamercolor{Feather}{fg=red!20,bg=red}

% Change the color of the structural elements:
%\setbeamercolor{structure}{fg=red}

% Change the frame title text color:
%\setbeamercolor{frametitle}{fg=blue}

% Change the normal text color background:
%\setbeamercolor{normal text}{fg=black,bg=gray!10}

%-------------------------------------------------------
% INCLUDE PACKAGES
%-------------------------------------------------------

\usepackage[utf8]{inputenc}
\usepackage[spanish,es-tabla]{babel}
\usepackage[T1]{fontenc}
\usepackage{helvet}
%-------------------------------------------------------
% DEFFINING AND REDEFINING COMMANDS
%-------------------------------------------------------

% colored hyperlinks
\newcommand{\chref}[2]{
  \href{#1}{{\usebeamercolor[bg]{Feather}#2}}
}

\makeatletter
\def\blfootnote{\gdef\@thefnmark{}\@footnotetext}
\makeatother

%-------------------------------------------------------
% INFORMATION IN THE TITLE PAGE
%-------------------------------------------------------

\title[The UNAM Feather Beamer Theme] % [] is optional - is placed on the bottom of the sidebar on every slide
{ % is placed on the title page
      \textbf{The UNAM Feather Beamer Theme}
}

\subtitle[The Feather Beamer Theme]
{
      \textbf{v. 1.0.0}
}

\author[Author Name]
{      \textbf{\large Author Name} \\
      {\ttfamily author.name@unam.mx}
}

\institute[UNAM]
{
      \textbf{Universidad Nacional Autónoma de México}\\
      Departamento de Física y Química Teórica\\
      Facultad de Química\\
      \vspace{0.2cm}
      \includegraphics[scale=0.15]{Feathergraphics/Ciencias-Quimicas.png}
      \vspace{-1cm}
  
  %there must be an empty line above this line - otherwise some unwanted space is added between the university and the country (I do not know why;( )
}

\date{\today}

%-------------------------------------------------------
% THE BODY OF THE PRESENTATION
%-------------------------------------------------------

\begin{document}

%-------------------------------------------------------
% THE TITLEPAGE
%-------------------------------------------------------

{\1% % this is the name of the PDF file for the background
\begin{frame}[plain,noframenumbering] % the plain option removes the header from the title page, noframenumbering removes the numbering of this frame only
  \titlepage % call the title page information from above
\end{frame}}

%-------------------------------------------------------
\section{Introducción}
%-------------------------------------------------------
\begin{frame}{Introducción}
%-------------------------------------------------------

  \begin{itemize}
    \item Esta es una presentación de beamer no oficial con el tema de colores de la UNAM. Se ha construido tomando como base el Template de Feather Beamer Theme creado por Lilyana Vankova bajo la licencia GNU General Public License v. 3 (GPLv3) http://www.gnu.org/licenses/.
    \item Para más información, consulte el Template original en: \url{https://www.overleaf.com/latex/templates/beamer-presentation-template-feather-theme/jcbpcdxqbxbf}
  \end{itemize}
\end{frame}

%-------------------------------------------------------
\section{Resultados}
%-------------------------------------------------------
\begin{frame}{Resultados 1}{Subtítulo Genérico}
%-------------------------------------------------------

  \begin{itemize}
    \item Se ha definido el color \textcolor{unamblue}{unamblue}
    \item Se ha definido el color \textcolor{unamgold}{unamgold}
  \end{itemize}

\blfootnote{Se puede usar blfootnote para referencias sin número.}

\end{frame}

\begin{frame}{Resultados 1}{Subtítulo Genérico}
%-------------------------------------------------------

\begin{block}{Bloque de texto}
  \begin{itemize}
    \item Información 1
    \item Información 2
    \item Información 3
    \item Información 4
  \end{itemize}
\end{block}

\blfootnote{Autor, D. Int. J. Quantum Power, 152(1), 125-130 (2023)}

\end{frame}

%-------------------------------------------------------
\section{Conclusiones}
%-------------------------------------------------------
\begin{frame}{Conclusiones}
%-------------------------------------------------------

  \begin{itemize}
    \item Los datos indican ...
    \item Se encontró una tendencia ...
  \end{itemize}

  \vspace{0.5cm}
  \textbf{\Large Perspectivas}
  \vspace{0.5cm}
  \begin{itemize}
    \item El camino a seguir es ...
    \item Se ha abierto el camino para lograr ...
  \end{itemize}

\end{frame}

\end{document}

\end{document}